\section{Related work}\label{sec:related}
\paragraph{Classical tour construction.} Insertion heuristics were analysed by Rosenkrantz, Stearns
and Lewis~\cite{rosenkrantz}. Their empirical behaviour on random uniform instances is catalogued by
Johnson and McGeoch~\cite{johnson} and the DIMACS TSP Challenge. Farthest insertion is the
strongest of the simple insertion rules. Savings~\cite{clarke} and Christofides are stronger at
large $n$. 2-opt~\cite{croes} and Or-opt~\cite{or} remove most of the remaining excess.
Rollout~\cite{bertsekas} is the classical way to turn a construction into a look-ahead policy, and
it is what several LLM-designed heuristics rediscover.

\paragraph{Neural and LLM-based heuristic design.} Kool et al.~\cite{kool} report insertion
heuristics next to learned constructive policies on exactly this distribution. We found no LLM-AHD paper that keeps that practice on the standard random test sets. Among the papers we found, only RouteRepair~\cite{routerepair}
places farthest insertion next to an LLM constructor: with its own adapter, its LLM construction is 9--11 points worse at $n=50$--$200$. It does not relate this to the MCTS-AHD leaderboard. RedAHD~\cite{redahd}
drops the per-step interface and writes whole-route functions with 2-opt. The critiques of LLM-AHD benchmarks we found target bin packing~\cite{herrmann,sim}, or called for stronger
\emph{search} baselines~\cite{zhang}, not classical algorithms. MOSAIC~\cite{mosaic} notes that
single-pass scoring prompts exclude look-ahead, and responds by permitting it.

\section{Discussion}\label{sec:discussion}
\paragraph{What the leaderboard measures.} With the whole matrix in every call, the step-by-step
interface constrains the \emph{output format}, not the algorithm. Any tour can be emitted, so the
ceiling is the optimum, and a short numpy function gets within a few percent of it. A published gap
on this benchmark therefore mixes two things: what the LLM search found, and how much computation
the found function is allowed per call. Recent gains come from the second. TIDE's heuristic runs
2-opt to convergence for up to 12 candidates at every step; HiFo-Prompt's runs a full
nearest-neighbour rollout per candidate, eight times. Without a compute budget, these numbers cannot
be compared with each other or with one-pass constructions.

\paragraph{Recommendations for this benchmark.}
\begin{enumerate}
\item \textbf{Report the classical baselines that fit the interface.} At minimum farthest
insertion, and a construction with 2-opt/Or-opt as the emitted reference.
\item \textbf{Report time per instance against those baselines on the same machine,} or fix a budget
per call.
\item \textbf{If the intent is to study one-pass construction rules, enforce it.} For example, pass
only the current node's row and the destination's row, or count matrix accesses.
\item \textbf{Publish per-instance reference values} with the test sets.
\end{enumerate}

\paragraph{Scope.} We study one benchmark task in one interface, on uniform instances up to 200
cities. We make no claim about LLM-AHD in general, or about frameworks in which the LLM designs
components of strong solvers (GLS, ACO, LKH parameters). Those tasks report strong solver baselines
already. Our timings come from a shared, loaded laptop; only orders of magnitude should be read
from them. Our point is narrower and, we think, actionable: on this widely used task, progress claims on random instances have been made against nearest neighbour, while textbook methods expressible inside the same interface were absent from the standard evaluation.
